\UseRawInputEncoding
\documentclass[12pt, openright, a4paper, openany, oneside, chapter=TITLE, section=TITLE, subsection=TITLE, subsubsection=TITLE, subsubsubsection=TITLE, subsubsubsection=TITLE, brazil]{abntex2}
%\usepackage{emoji}
\usepackage{Pacotes} % Pacote mantido do seu template original
\usepackage{lipsum} % Pacote mantido do seu template original
\usepackage{dirtytalk} % Pacote mantido do seu template original
\usepackage{listings} % Pacote mantido do seu template original
\usepackage{xcolor} % Pacote mantido do seu template original

% --- Definições de cores e estilos mantidas do seu template ---
\definecolor{backcolour}{rgb}{0.97,0.97,0.97} % fundo branco
\definecolor{codegreen}{rgb}{0.25,0.5,0.25}    % comentários (verde acinzentado)
\definecolor{codegray}{rgb}{0.4,0.4,0.4}       % textos auxiliares
\definecolor{codepurple}{rgb}{0.0,0.0,1.0}     % palavras-chave (azul)
\definecolor{codeorange}{rgb}{0.8,0.3,0.0}     % funções (laranja)
\definecolor{black}{rgb}{0,0,0}               % texto padrão

\lstdefinestyle{mystyle}{
    backgroundcolor=\color{backcolour},   
    commentstyle=\color{codegreen},
    keywordstyle=\color{codepurple}\bfseries,
    stringstyle=\color{codeorange},
    numberstyle=\tiny\color{codegray},
    basicstyle=\ttfamily\color{black},
    breakatwhitespace=false,         
    breaklines=true,                 
    captionpos=b,                    
    keepspaces=true,                 
    numbers=left,                    
    numbersep=5pt,                   
    showspaces=false,                
    showstringspaces=false,
    showtabs=false,                  
    tabsize=2
}
\lstset{style=mystyle}

%\usepackage{showframe}
\everymath{\displaystyle}
\checkandfixthelayout

% --- METADADOS DO RELATÓRIO DE AUDITORIA ---
\titulo{Relatório de Auditoria de Produto: Módulo de Interruptor Capacitivo Bi-Estável}
\autor{Denis Natã Henden Loyola (GRR20231589)\\ Ricardo Marin Israel (GRR20231588)}
\local{Curitiba}
\data{2025}
\instituicao{Universidade Federal do Paraná}
\tipotrabalho{Relatório de Auditoria} % Tipo de trabalho atualizado
\preambulo{Relatório de auditoria de produto apresentado à disciplina TE366 - Desenho Universal, do curso de Engenharia Elétrica, Setor de Tecnologia, Universidade Federal do Paraná (UFPR), como requisito parcial para avaliação do Trabalho 02. \\
Professor responsável: Prof. Dr. Sebastião Ribeiro Júnior.}
% --- FIM DOS METADADOS ---

\makeatletter
\hypersetup{
pdftitle={\@title},
pdfauthor={\@author},
pdfsubject={\imprimirpreambulo},
pdfkeywords={Auditoria}{Desenho Universal}{Acessibilidade}{Interruptor Capacitivo}{TE366}{NBR 9050}, % Palavras-chave atualizadas
pdfcreator={LaTeX with abnTeX2},
colorlinks=true,
linkcolor=black,
citecolor=black,
urlcolor=gray,
}
\makeatother


\begin{document}
\pretextual

\imprimircapa
\imprimirfolhaderosto
\clearpage

% --- RESUMO DO RELATÓRIO DE AUDITORIA ---
\chapter*{Resumo}
Este relatório de auditoria técnica analisa o projeto "Módulo de Interruptor Capacitivo Bi-Estável", desenvolvido para a disciplina TE366. O objetivo é verificar a conformidade do projeto e sua documentação (Memorial Descritivo e Desenho Técnico) com os 7 Princípios do Desenho Universal e as normas ABNT NBR 9050 e NBR 5410. A análise foca na solução de acessibilidade (interface capacitiva) e na correção da documentação de engenharia. A auditoria identificou pontos de conformidade, como a aplicação do princípio de "Baixo Esforço Físico", mas também apontou não conformidades. Os principais problemas são a ausência de feedback sensorial (Princípio "Informação Perceptível") e falhas na padronização do desenho técnico (cotagem e representação de vistas). O relatório termina com recomendações para adequar o projeto às normas.
% --- FIM DO RESUMO ---

\clearpage
\tableofcontents*
\cleardoublepage

\textual

% --- CONTEÚDO DO RELATÓRIO DE AUDITORIA ---

\chapter{Informações Gerais}
\label{sec:info_gerais}

Seguindo o modelo de auditoria da Aula 09 da disciplina TE366, este relatório detalha a análise do projeto "Módulo de Interruptor Capacitivo Bi-Estável". A avaliação é baseada em critérios técnicos e normas. A auditoria visa garantir que o projeto atenda aos requisitos de acessibilidade e segurança, conforme os objetivos da disciplina.

\begin{itemize}
    \item \textbf{Data da Auditoria:} 31 de outubro de 2025.
    \item \textbf{Local:} Universidade Federal do Paraná - Setor de Tecnologia (Auditoria baseada em revisão documental).
    \item \textbf{Auditor:} Denis Natã Henden Loyola (GRR20231589) e Ricardo Marin Israel (GRR20231588).
    \item \textbf{Projeto Auditado (Autores):} Denis Natã Henden Loyola (GRR20231589) e Ricardo Marin Israel (GRR20231588).
    \item \textbf{Objetivo da Auditoria:} Verificar a conformidade do projeto com os requisitos de acessibilidade da NBR 9050, os padrões de segurança elétrica da NBR 5410, e as normas de desenho técnico aplicáveis.
\end{itemize}


\chapter{Escopo da Auditoria}
\label{sec:escopo}

A auditoria cobriu os seguintes elementos do projeto, conforme o Memorial Descritivo fornecido pelos autores:

\begin{itemize}
    \item Concepção estrutural e dimensional, e a compatibilidade com padrões de mercado.
    \item Especificação de materiais do invólucro e superfície de acionamento, avaliando propriedades funcionais e de segurança.
    \item Análise da aplicação dos 7 Princípios de Desenho Universal no produto.
    \item Princípio de funcionamento elétrico, premissas de projeto e limitações operacionais.
    \item Análise das Vistas Ortográficas (Apêndice A) e sua conformidade com normas de desenho técnico (cotagem, simbologia, legendas).
\end{itemize}


\chapter{Normas Técnicas de Referência}
\label{sec:normas}

Esta auditoria baseou-se nas seguintes normas técnicas, referenciadas no projeto e no material da disciplina TE366:

\begin{itemize}
    \item \textbf{ABNT NBR 9050:} Norma que estabelece critérios e parâmetros técnicos de acessibilidade a edificações, mobiliário, espaços e equipamentos urbanos.
    \item \textbf{ABNT NBR 5410:} Norma que fixa as condições para instalações elétricas de baixa tensão, visando a segurança de pessoas e o funcionamento adequado da instalação.
    \item \textbf{ABNT NBR 8196:} Norma para emprego de escalas em desenho técnico.
    \item \textbf{ABNT NBR 8403:} Norma para aplicação de linhas em desenhos técnicos (tipos e larguras).
    \item \textbf{Princípios do Desenho Universal:} Conjunto de sete princípios conforme apresentado na Aula 01 da disciplina (Uso Equitativo, Flexibilidade no Uso, Uso Simples e Intuitivo, Informação Perceptível, Tolerância ao Erro, Baixo Esforço Físico e Tamanho e Espaço para Aproximação e Uso).
\end{itemize}


\chapter{Metodologia da Auditoria}
\label{sec:metodologia}

A auditoria foi realizada em uma fase de investigação:

\begin{itemize}
    \item \textbf{Revisão Documental:} Verificação do "Memorial Descritivo – Módulo de Interruptor Capacitivo Bi-Estável" e do Apêndice A (Vistas Ortográficas), comparando as soluções propostas com os requisitos das normas citadas. A análise incluiu a identificação de pontos de conformidade e não conformidade.
\end{itemize}


\chapter{Resultados da Auditoria}
\label{sec:resultados}

A análise técnica identificou os seguintes pontos de conformidade e não conformidade, detalhados abaixo.

\section{Aplicação dos Princípios de Desenho Universal}

\subsection{Conformidade}
O projeto adere a princípios do Desenho Universal, buscando ser uma alternativa acessível para usuários com limitações motoras. A conformidade é vista nos seguintes aspectos:

\begin{itemize}
    \item \textbf{Uso Equitativo (Princípio 1) e Baixo Esforço Físico (Princípio 6):} A substituição do acionamento mecânico por uma interface capacitiva elimina a necessidade de força pontual ou precisão motora. A "área de acionamento" (toda a face frontal) permite a comutação por um simples contato, que pode ser feito com cotovelos, punhos ou outras partes do corpo, o que representa uma melhoria na acessibilidade motora.
    
    \item \textbf{Facilidade de Higienização (Derivado do Princípio 5 - Tolerância ao Erro):} A superfície frontal selada, sem frestas ou reentrâncias, facilita a limpeza. Isso é útil para todos os usuários e importante em ambientes como hospitais ou laboratórios.
\end{itemize}

\subsection{Não Conformidades}
\begin{itemize}
    \item \textbf{Informação Perceptível (Princípio 4):} O memorial descritivo não detalha mecanismos de feedback sensorial para o usuário, o que é uma lacuna no projeto:
        \begin{itemize}
            \item \textit{Não Conformidade 5.1.1:} Ausência de um indicador de estado claro (ligado/desligado). Um indicador visual (como um LED de alto contraste) ou tátil é fundamental. Sem ele, usuários com deficiência visual ou auditiva (em caso de lâmpadas silenciosas) não conseguem verificar se o circuito está ativo ou inativo, o que anula parte do benefício de acessibilidade.
            \item \textit{Não Conformidade 5.1.2:} Ausência de um feedback de acionamento imediato (sonoro ou tátil) que confirme o registro do toque. Em uma superfície capacitiva sem partes móveis, o usuário (especialmente um com deficiência visual ou neuropatia) pode tocar a superfície e não ter certeza se o comando foi efetivamente aceito e processado pelo sistema, gerando insegurança.
        \end{itemize}
\end{itemize}

\section{Especificações Técnicas e Segurança (NBR 5410)}

\subsection{Conformidade}
\begin{itemize}
    \item O projeto prevê o uso de invólucro em termoplástico (PVC) com propriedade antichama e isolamento galvânico do circuito de controle (optoacoplador), atendendo às premissas de segurança contra choque elétrico da NBR 5410.
    \item O memorial descreve as limitações de aplicação (cargas indutivas) e as premissas de instalação para o funcionamento seguro.
\end{itemize}

\subsection{Não Conformidades}
\begin{itemize}
    \item \textit{Não Conformidade 5.2.1:} A dependência do condutor Neutro no ponto de instalação é uma barreira técnica significativa. O projeto assume uma instalação moderna (padrão Fase-Neutro-Retorno na caixa $4 \times 2"$). Isso limita a aplicação em "instalações elétricas prediais existentes", visto que muitas edificações mais antigas (especialmente antes da NBR 5410 de 1997) utilizam apenas Fase e Retorno na caixa do interruptor, inviabilizando a alimentação do circuito eletrônico do dispositivo.
\end{itemize}

\section{Desenho Técnico (Apêndice A do Memorial)}

Avaliação técnica baseada nos critérios de auditoria para Desenho Técnico (Aula 09).

\subsection{Conformidade}
\begin{itemize}
    \item O desenho apresenta as vistas ortográficas principais (Superior, Frontal e Laterais) em primeiro diedro, conforme indicado na Figura 1 do memorial, permitindo a compreensão geométrica básica.
\end{itemize}

\subsection{Não Conformidades}
\begin{itemize}
    \item \textit{Não Conformidade 5.3.1 (Sistema de Cotagem):} O sistema de cotagem apresenta problemas de clareza. Há sobreposição de linhas auxiliares e valores dimensionais, dificultando a legibilidade. Além disso, o uso de precisão excessiva (ex: 6,91 mm; 6,71 mm; R1,88 mm) não parece justificado para a aplicação (um invólucro de interruptor) e pode implicar em tolerâncias de fabricação desnecessariamente restritas e caras, sem agregar valor funcional ao produto.
    
    \item \textit{Não Conformidade 5.3.2 (Simbologia e Representação Gráfica):} Foi identificada uma mistura de convenções de representação. As vistas laterais e superior exibem detalhes internos (como bornes de conexão e furos de fixação da PCI) com linhas contínuas, como se fossem elementos visíveis externamente. A representação correta exigiria o uso de linhas tracejadas (para arestas não visíveis, conforme NBR 8403) ou a aplicação de vistas de corte (seção) para expor a montagem interna. A forma atual compromete a interpretação da geometria e da montagem do produto.
\end{itemize}


\chapter{Recomendações}
\label{sec:recomendacoes}

Com base nas não conformidades identificadas na Seção \ref{sec:resultados}, recomendam-se as seguintes ações corretivas:

\begin{enumerate}
    \item \textbf{Revisão de Projeto (Feedback):} Incluir obrigatoriamente um LED indicador de estado (Ligado/Desligado) de fácil visualização e um sistema de feedback de acionamento (sonoro, como um buzzer, ou tátil, como um micro-vibrador) para confirmação inequívoca do toque. Isso é essencial para atender ao Princípio 4 (Informação Perceptível).
    
    \item \textbf{Revisão de Projeto (Compatibilidade):} Avaliar a viabilidade técnica de uma versão do circuito que não dependa do condutor Neutro. Isso pode envolver, por exemplo, um circuito de alimentação "parasita" que utiliza a própria corrente da carga (quando esta permite uma pequena corrente de fuga) ou uma solução com bateria, ampliando a compatibilidade com instalações elétricas mais antigas.
    
    \item \textbf{Correção do Desenho (Cotagem):} Revisar o sistema de cotagem no Apêndice A. As cotas devem ser reorganizadas para eliminar sobreposições de linhas. Recomenda-se também a padronização da precisão dimensional (número de casas decimais) para que seja condizente com o processo de fabricação do invólucro (ex: injeção plástica), evitando tolerâncias desnecessariamente restritas.
    
    \item \textbf{Correção do Desenho (Convenções):} Aplicar corretamente as convenções de desenho técnico. Componentes internos (bornes, PCI) devem ser representados com linhas tracejadas (NBR 8403) quando ocultos, ou deve-se utilizar vistas de corte (seção) para detalhar a montagem interna, separando claramente o que é externo do que é interno.
\end{enumerate}


\chapter{Conclusão}
\label{sec:conclusao}

O projeto do "Módulo de Interruptor Capacitivo Bi-Estável" atende aos requisitos de baixo esforço físico e uso equitativo, sendo uma solução de acessibilidade viável para usuários com limitações motoras.
\par
Apesar dos pontos positivos, a auditoria identificou a necessidade de ajustes para garantir a acessibilidade (principalmente com a implementação de feedback perceptível) e para adequar o desenho técnico aos padrões normativos da engenharia, especialmente na clareza da representação de cotas e vistas.
\par
\vspace{0.5cm}
\noindent \textbf{Status Final da Auditoria: Aprovado com Ressalvas.}
\par
\vspace{0.2cm}
\noindent Recomenda-se que as implementações listadas na Seção \ref{sec:recomendacoes} sejam realizadas e testadas antes da reapresentação do projeto para a validação final do Trabalho 02 da disciplina TE366.

% --- BIBLIOGRAFIA ---
\begin{thebibliography}{99}

\bibitem{nbr9050}
ASSOCIAÇÃO BRASILEIRA DE NORMAS TÉCNICAS. \textit{NBR 9050: Acessibilidade a edificações, mobiliário, espaços e equipamentos urbanos}. Rio de Janeiro, 2020.

\bibitem{nbr5410}
ASSOCIAÇÃO BRASILEIRA DE NORMAS TÉCNICAS. \textit{NBR 5410: Instalações elétricas de baixa tensão}. Rio de Janeiro, 2004.

\bibitem{nbr8196}
ASSOCIAÇÃO BRASILEIRA DE NORMAS TÉCNICAS. \textit{NBR 8196: Desenho técnico — Emprego de escalas}. Rio de Janeiro, 1999.

\bibitem{nbr8403}
ASSOCIAÇÃO BRASILEIRA DE NORMAS TÉCNICAS. \textit{NBR 8403: Desenho técnico — Aplicação de linhas em desenhos — Tipos de linhas — Larguras das linhas}. Rio de Janeiro, 1984.

\end{thebibliography}

% --- APÊNDICES (Removidos, pois não fazem parte deste relatório) ---

\end{document}